\begin{register}{H}{mhwloop0\_start\_addr}{0x7C6}\label{regdesc:MHWLOOP0_START_ADDR}
\regfield{SA}{32}{0}{{0x00000000}}%
\reglabel{Reset}\regnewline%
\begin{regdesc}\begin{reglist}
\label{fielddesc:HWLP0_SA}\item [SA] 配置 HW Loop0 第一条指令的 32 位地址。当 \hyperref[fielddesc:HWSTATE]{mhwloop\_state\_reg.STATE} 不为 0x0 时有效，即不为 OFF。(R/W)
\end{reglist}\end{regdesc}
\end{register}

\begin{register}{H}{mhwLoop0 \_end\_addr}{0x7C7}\label{regdesc:MHWLOOP0_END_ADDR}
\regfield{EA}{32}{0}{{0x00000000}}%
\reglabel{Reset}\regnewline%
\begin{regdesc}\begin{reglist}
\label{fielddesc:HWLP0_EA}\item [EA] 配置 HW Loop0 最后一条指令的 32 位地址。当 \hyperref[fielddesc:HWSTATE]{mhwloop\_state\_reg.STATE} 不为 0x0 时有效，即不为 OFF。(R/W)
\end{reglist}\end{regdesc}
\end{register}

\begin{register}{H}{mhwloop0\_count}{0x7C8}\label{regdesc:MHWLOOP0_COUNT}
\regfield{(reserved)}{20}{0}{{0x00000}}%
\regfield{COUNT}{12}{0}{{0x000}}%
\reglabel{Reset}\regnewline%
\begin{regdesc}\begin{reglist}
\label{fielddesc:HWLP0_CNT}\item [COUNT] 配置 HW Loop0 的循环次数。当 \hyperref[fielddesc:HWSTATE]{mhwloop\_state\_reg.STATE} 不为 0x0 时有效，即不为 OFF。(R/W)
\end{reglist}\end{regdesc}
\end{register}

\begin{register}{H}{mhwloop1\_start\_addr}{0x7C9}\label{regdesc:MHWLOOP1_START_ADDR}
\regfield{SA}{32}{0}{{0x00000000}}%
\reglabel{Reset}\regnewline%
\begin{regdesc}\begin{reglist}
\label{fielddesc:HWLP1_SA}\item [SA] 配置 HW Loop1 第一条指令的 32 位地址。当 \hyperref[fielddesc:HWSTATE]{mhwloop\_state\_reg.STATE} 不为 0x0 时有效，即不为 OFF。(R/W)
\end{reglist}\end{regdesc}
\end{register}

\begin{register}{H}{mhwloop1\_end\_addr}{0x7CA}\label{regdesc:mhloop1_end_addr}
\regfield{EA}{32}{0}{{0x00000000}}%
\reglabel{Reset}\regnewline%
\begin{regdesc}\begin{reglist}
\label{fielddesc:HWLP1_EA}\item [EA] 配置 HW Loop1 最后一条指令的 32 位地址。当 \hyperref[fielddesc:HWSTATE]{mhwloop\_state\_reg.STATE} 不为 0x0 时有效，即不为 OFF。(R/W)
\end{reglist}\end{regdesc}
\end{register}

\begin{register}{H}{mhwloop1\_count}{0x7CB}\label{regdesc:mhloop1_count}
\regfield{(reserved)}{20}{0}{{0x00000}}%
\regfield{COUNT}{12}{0}{{0x000}}%
\reglabel{Reset}\regnewline%
\begin{regdesc}\begin{reglist}
\label{fielddesc:MHLP1_COUNT}\item [COUNT] 配置 HW Loop1 的循环次数。当mhwloop\_state\_reg.STATE 不为 0x0 时有效，即不为 OFF。(R/W)
\end{reglist}\end{regdesc}
\end{register}

\begin{register}{H}{mext\_ill\_reg}{0x7F0}\label{regdesc:mext_ill_reg}
\regfield{(reserved)}{29}{3}{{0x000000}}%
\regfield{ILL\_PIE}{1}{2}{{0x0}}%
\regfield{ILL\_HWL}{1}{1}{{0x0}}%
\regfield{ILL\_FPU}{1}{0}{{0x0}}%
\reglabel{Reset}\regnewline%
\begin{regdesc}\begin{reglist}
\label{fielddesc:ILL_FPU}\item [ILL\_FPU] 表示是否由于禁用 F 扩展而触发非法指令。\\
0：未触发\\
1：已触发\\
(RO)
\label{fielddesc:ILL_HWL}\item [ILL\_HWL] 表示是否由于禁用 HWLP 扩展而触发非法指令。\\
0：未触发\\
1：已触发\\
(RO)
\label{fielddesc:ILL_PIE}\item [ILL\_PIE] 表示是否由于禁用 PIE 扩展而触发非法指令。\\
0：未触发\\
1：已触发\\
(RO)
\end{reglist}\end{regdesc}
\end{register}

\begin{register}{H}{mhwloop\_state\_reg}{0x7F1}\label{regdesc:mhwlp_state_reg}
\regfield{(reserved)}{29}{3}{{0x000000}}%
\regfield{URW}{1}{2}{{0x0}}%
\regfield{STATE}{2}{0}{{0x0}}%
\reglabel{Reset}\regnewline%
\begin{regdesc}\begin{reglist}
\label{fielddesc:HWSTATE}\item [STATE] 配置 HWLP 扩展的状态。\\
0x0：关闭（访问 HWLP 指令和 CSR 将触发非法指令异常）\\
0x1：初始\\
0x2：干净\\
0x3：脏\\
(R/W)
\label{fielddesc:URW}\item [URW] 配置用户模式下对用户模式 HWLOOP CSR 的访问。\\
0：禁用用户模式访问\\
1：启用用户模式访问\\
(R/W)
\end{reglist}\end{regdesc}
\end{register}

\begin{register}{H}{uhwloop0\_start\_addr}{0x8C0}\label{regdesc:uhloop0_start_addr}
\regfield{SA}{32}{0}{{0x00000000}}%
\reglabel{Reset}\regnewline%
\begin{regdesc}\begin{reglist}
\label{fielddesc:UHWLP0_SA}\item [SA] 配置用户模式下 HW Loop0 第一条指令的 32 位地址。当 \hyperref[fielddesc:HWSTATE]{mhwloop\_state\_reg.STATE} 不为 0x0 时有效，即不为 OFF。仅当 \hyperref[fielddesc:HWSTATE]{mhwloop\_state\_reg.URW} 为 1 时可读写。(R/W)
\end{reglist}\end{regdesc}
\end{register}

\begin{register}{H}{uhwloop0\_end\_addr}{0x8C1}\label{regdesc:uhloop0_end_addr}
\regfield{EA}{32}{0}{{0x00000000}}%
\reglabel{Reset}\regnewline%
\begin{regdesc}\begin{reglist}
\label{fielddesc:UHWLP0_EA}\item [EA] 配置用户模式下 HW Loop0 最后一条指令的 32 位地址。当 \hyperref[fielddesc:HWSTATE]{mhwloop\_state\_reg.STATE} 不为 0x0 时有效，即不为 OFF。仅当 \hyperref[fielddesc:HWSTATE]{mhwloop\_state\_reg.URW} 为 1 时可读写。(R/W)
\end{reglist}\end{regdesc}
\end{register}

\begin{register}{H}{uhwloop0\_count}{0x8C2}\label{regdesc:uhloop0_count}
\regfield{(reserved)}{20}{0}{{0x00000}}%
\regfield{Count}{12}{0}{{0x000}}%
\reglabel{Reset}\regnewline%
\begin{regdesc}\begin{reglist}
\label{fielddesc:UHWLP0_CNT}\item [Count] 配置用户模式下 HW Loop0 的循环次数。当 \hyperref[fielddesc:HWSTATE]{mhwloop\_state\_reg.STATE} 不为 0x0 时有效，即不为 OFF。仅当 \hyperref[fielddesc:HWSTATE]{mhwloop\_state\_reg.URW} 为 1 时可读写。(R/W)
\end{reglist}\end{regdesc}
\end{register}

\begin{register}{H}{uhwloop1\_start\_addr}{0x8C3}\label{regdesc:uhloop1_start_addr}
\regfield{SA}{32}{0}{{0x00000000}}%
\reglabel{Reset}\regnewline%
\begin{regdesc}\begin{reglist}
\label{fielddesc:UHWLP1_SA}\item [SA] 配置用户模式下 HW Loop1 的第一条指令的 32 位地址。当 \hyperref[fielddesc:HWSTATE]{mhwloop\_state\_reg.STATE} 不为 0x0 时有效，即不为 OFF。仅当 \hyperref[fielddesc:HWSTATE]{mhwloop\_state\_reg.URW} 为 1 时可读写。(R/W)
\end{reglist}\end{regdesc}
\end{register}

\begin{register}{H}{uhwloop1\_end\_addr}{0x8C4}\label{regdesc:uhloop1_end_addr}
\regfield{EA}{32}{0}{{0x00000000}}%
\reglabel{Reset}\regnewline%
\begin{regdesc}\begin{reglist}
\label{fielddesc:UHWLP1_EA}\item [EA] 配置用户模式下 HW Loop1 的最后一条指令的 32 位地址。当 \hyperref[fielddesc:HWSTATE]{mhwloop\_state\_reg.STATE} 不为 0x0 时有效，即不为 OFF。仅当 \hyperref[fielddesc:HWSTATE]{mhwloop\_state\_reg.URW} 为 1 时可读写。(R/W)
\end{reglist}\end{regdesc}
\end{register}

\begin{register}{H}{uhwloop1\_count}{0x8C5}\label{regdesc:uhloop1_count}
\regfield{(reserved)}{20}{0}{{0x00000}}%
\regfield{Count}{12}{0}{{0x000}}%
\reglabel{Reset}\regnewline%
\begin{regdesc}\begin{reglist}
\label{fielddesc:UHWLP1_CNT}\item [Count] 配置用户模式下 HW Loop1 的循环次数。当 \hyperref[fielddesc:HWSTATE]{mhwloop\_state\_reg.STATE} 不为 0x0 时有效，即不为 OFF。仅当 \hyperref[fielddesc:HWSTATE]{mhwloop\_state\_reg.URW} 为 1 时可读写。(R/W)
\end{reglist}\end{regdesc}
\end{register}